\begin{myChapterIntro}{本章涵盖}
\begin{itemize}
\item 单例（Singleton）设计模式
\item 组合（Composite）设计模式
\item 装饰器（Decorator）设计模式
\end{itemize}
\end{myChapterIntro}

本章的设计模式为求解架构问题提供了模型：在这类问题中，我们必须管理一个由一或多个对象组成的集合，使其以统一的方式运作。应用恰当的设计模式可以简化操作这些对象的应用代码，还能让代码更灵活。

\myGraphic{0.893}{images/CH14_00_UN01_Mak.png}{}

单例设计模式为最简单的情形——集合中只有一个对象——建模。在应用的运行期间，其类的实例最多只能存在一个，即单例对象。

组合设计模式为管理存储在层级树结构中的对象的应用提供了模型。如果应用能够像对待对象组合那样对待单个对象，就能大幅降低其代码的复杂度。

我们可能希望应用在运行时以属性和行为的形式为对象添加职责，却无需修改该对象的代码。装饰器设计模式把这些以对象形式存在的额外职责称为“装饰”（decorations），并提供了一个灵活处理对象装饰的模型。

\begin{myWarning}{注意}
请务必阅读第 4 部分的引言，了解关于设计模式整体的重要信息，并弄清楚本章及后续各章是如何讲解每一种设计模式的。
\end{myWarning}

\section{单例设计模式确保一个类只有一个对象}

假设你的应用需要向一台或多台打印机输出打印内容。每当应用需要打印时，它都必须访问一个 \texttt{PrintSpooler} 对象。应用并不知道这个后台处理程序（spooler）对象是否已经存在。如果该对象不存在，就会创建一个，而应用对此毫不知情。每当应用要打印时，它访问的都是同一个 \texttt{PrintSpooler} 对象，因此至多只存在一个 \texttt{PrintSpooler} 对象。所以，该对象就是一个\emph{单例}对象。单例设计模式为支持此类对象的软件架构提供了模型。

在具体示例中，假设大学体育馆有一个贵宾包厢（executive suite），仅在特定的体育赛事期间可供使用。只有一张通行证，允许持有人在包厢可用时使用它。这张通行证就是一个单例对象，而且至多只能有一张。

\subsection{14.1.1 期望的设计特性}

一个至多只能拥有某个特定类的一个实例的应用，应当具备以下设计特性：

\begin{itemize}
\item \emph{DF 1}——应用必须有一种方式来访问单例对象。
\item \emph{DF 2}——无论何时应用需要单例对象，它访问的都必须是同一个对象。
\item \emph{DF 3}——如果单例对象尚未存在，必须在应用不采取任何操作的情况下创建一个。
\item \emph{DF 4}——如果应用从不访问该对象，那么单例对象无需存在。
\item \emph{DF 5}——单例对象的实例永远不能超过一个。
\end{itemize}

\subsection{14.1.2 使用单例设计模式之前}

我们应用的第一个版本不使用单例设计模式，因此无法具备全部期望的设计特性。随后，我们将使用该设计模式来重构这个应用。

实现单例对象的一个显而易见的方法，是把它作为全局变量的值。在我们应用的第一个版本中，类 \texttt{ExecutivePass} 就是贵宾通行证的一个直白实现（图 14.1）。

\myGraphic{0.266}{images/CH14_F01_Mak.png}{图 14.1 我们对单例类的首次尝试}

成员函数 \texttt{obtain()} 允许某个具名用户获得并持有这张通行证（代码清单 14.1）。这张通行证有一个随机生成的四位 \texttt{key} 用于标识它。成员函数 \texttt{print()} 打印 \texttt{holder} 的名字和 \texttt{key} 的值。

就像我们在 6.2.2 节对 \texttt{Date} 类所做的那样，我们将重写并插桩拷贝构造函数和拷贝赋值运算符。这样做将展示：如果单例对象设计不当，它是如何失效的。

\begin{codeboxt}{代码清单 14.1（程序 14.1 ExecPass）：\texttt{ExecutivePass.h}（使用设计模式之前）}
#include <iostream>
#include <string>
#include <stdlib.h>

using namespace std;

class ExecutivePass
{
public:
    ExecutivePass() : key(generate_key()) {}                     ①

    ExecutivePass(ExecutivePass& other) : key(generate_key()) {} ②

    ExecutivePass& operator =(ExecutivePass& other)              ③
    {
        key = generate_key();
        return *this;
    }

    void obtain(const string hldr) { holder = hldr; }

    void print()
    {
        cout << "Executive pass now held by " << holder
             << ", key = " << key << endl;
    }

private:
    int key;
    string holder;

    int generate_key() const { return rand()%10000 + 1000; }
};
\end{codeboxt}
\noindent{\fontsize{9bp}{12bp}\selectfont ① 默认构造函数}\par
\noindent{\fontsize{9bp}{12bp}\selectfont ② 拷贝构造函数}\par
\noindent{\fontsize{9bp}{12bp}\selectfont ③ 拷贝赋值运算符}\par

在测试程序中，全局变量 \texttt{pass} 的值就是那个 \texttt{ExecutivePass} 对象。

\begin{codeboxt}{代码清单 14.2（程序 14.1 ExecPass）：\texttt{tester.cpp}（使用设计模式之前）}
#include <iostream>
#include "ExecutivePass.h"

ExecutivePass pass;    

using namespace std;

int main()
{
    pass.obtain("Ron");             ①
    cout << "obtain: "; pass.print();

    pass.obtain("Sal");             ②
    cout << "obtain: "; pass.print();

    ExecutivePass copy1;            ③
    copy1 = pass;                   ③
    copy1.obtain("Bob");            ③

    cout << "copy 1: "; copy1.print();
    ExecutivePass copy2(pass);      ④
    copy2.obtain("Flo");            ④
    cout << "copy 2: "; copy2.print();

    return 0;
}
\end{codeboxt}
\noindent{\fontsize{9bp}{12bp}\selectfont ① Ron 获得全局通行证。}\par
\noindent{\fontsize{9bp}{12bp}\selectfont ② 把全局通行证转交给 Sal。}\par
\noindent{\fontsize{9bp}{12bp}\selectfont ③ 作为赋值目标的通行证本地副本}\par
\noindent{\fontsize{9bp}{12bp}\selectfont ④ 由拷贝构造函数创建的通行证本地副本}\par

程序的输出如下

\begin{codebox}
obtain: Executive pass now held by Ron, key = 7807
obtain: Executive pass now held by Sal, key = 7807
copy 1: Executive pass now held by Bob, key = 1073
copy 2: Executive pass now held by Flo, key = 4658
\end{codebox}

显然，我们应用的这一版本失败了。用户 Ron 获取了带有其键值的全局 \texttt{pass} 对象。同一个对象被成功转交给了 Sal。但我们却能够通过赋值创建通行证的一个副本 \texttt{copy1}，并被用户 Bob 获取，它有自己的键值。最后，我们又用拷贝构造函数创建了 \texttt{copy2}，被用户 Flo 获取，它也有自己的键值。最终，这个本应是单例的对象有了三个副本。

\myGraphic{0.825}{images/CH14_F01_UN02_Mak.png}{}

我们应用第一个版本的主要缺陷有

\begin{itemize}
\item \emph{使用全局变量}——在设计良好的应用中不应使用全局变量。应用的不同部分都能访问全局变量。某一部分可能改变变量的值，而其他并未预料到这一改变的部分却毫不知情，从而导致意外的麻烦和可能的逻辑错误。
\item \emph{创建顺序}——如果声明在不同源文件中的不同全局对象相互依赖，就可能发生逻辑错误，因为应用启动时这些对象的创建顺序是未定义的。
\item \emph{对象始终存在}——在我们的示例中，贵宾包厢并非每场体育赛事都开放。由于我们把通行证定义成了全局的，它就始终存在。如果单例对象构造代价高昂或消耗大量资源，这就会成为问题。而且我们无法删除这个全局对象。
\item \emph{本地副本}——我们能够为单例对象创建额外的本地副本。
\item \emph{可以拷贝和赋值该对象}——这些操作会为这个本应是单例的对象创建更多实例。
\end{itemize}

\subsection{14.1.3 使用单例设计模式之后}

我们可以通过重构代码、使其仿照单例设计模式来修复类 \texttt{ExecutivePass} 存在的问题。

\begin{mySidebar}{单例设计模式}

“确保一个类只有一个实例，并提供一个访问它的全局访问点”（GoF 127）。
\end{mySidebar}

我们必须做出几处重大改动（图 14.2）：

\begin{itemize}
\item 我们必须把构造函数设为私有。不能允许代码声明一个 \texttt{ExecutivePass} 对象，也不能用 \texttt{new} 动态创建它。我们必须控制单例通行证对象的创建，并且只应在需要时才创建单例对象（惰性构造）。
\item 我们必须禁用拷贝构造函数和拷贝赋值运算符。不能允许代码为单例对象创建副本。
\item 一个私有的静态 \texttt{ExecutivePass} 成员变量必须指向那个唯一的单例对象（如果它存在的话）。否则，该变量的值必须是 \texttt{nullptr}。
\item 一个公有的静态 \texttt{obtain()} 成员函数必须返回指向单例对象的指针，并设置谁持有这张通行证。如果对象尚未存在，它必须创建该对象。\texttt{ExecutivePass::obtain()} 就是访问单例对象的全局访问点。
\item 删除单例对象之后，必须能够稍后再创建一个新的。
\end{itemize}

\myGraphic{0.793}{images/CH14_F02_Mak.png}{图 14.2 一个设计良好的单例对象。构造函数是私有的。拷贝构造函数和拷贝赋值运算符都被禁用。私有静态成员变量 \texttt{single\_instance} 指向该单例对象，公有静态成员函数 \texttt{obtain()} 返回该指针的值。}

我们把构造函数设为私有，并禁用拷贝构造函数和拷贝赋值运算符。

\begin{codeboxt}{代码清单 14.3（程序 14.2 ExecPass-SingletonDP）：\texttt{ExecutivePass.h}}
#include <iostream>
#include <string>
#include <stdlib.h>

using namespace std;

class ExecutivePass
{
public:
    ~ExecutivePass();

    ExecutivePass(ExecutivePass& other) = delete;              ①
    ExecutivePass& operator =(ExecutivePass& other) = delete;  ①

    int get_key() const { return key; }

    static ExecutivePass *obtain(const string holder);

    void print()
    {
        cout << "Executive pass now held by " << holder
             << ", key = " << key << endl;
    }

private:
    ExecutivePass() : key(generate_key()) {}                   ②

    static ExecutivePass *single_instance;                     ③

    int key;
    string holder;

    int generate_key() const { return rand()%10000 + 1000; }
};
\end{codeboxt}
\noindent{\fontsize{9bp}{12bp}\selectfont ① 禁用拷贝构造函数和拷贝赋值运算符。}\par
\noindent{\fontsize{9bp}{12bp}\selectfont ② 私有构造函数}\par
\noindent{\fontsize{9bp}{12bp}\selectfont ③ 指向单例对象的静态指针}\par

私有静态成员变量 \texttt{single\_instance} 指向单例对象（如果它存在的话）；否则其值为 \texttt{nullptr}（代码清单 14.4）。公有静态成员函数 \texttt{obtain()} 返回指向已存在单例对象的指针。如果单例对象尚未存在，该函数会创建并返回指向一个新创建单例对象的指针。析构函数把 \texttt{single\_instance} 设为 \texttt{nullptr}，这样在应用运行期间，当下次我们的代码需要单例对象时，就能创建一个新的。

\begin{codeboxt}{代码清单 14.4（程序 14.2 ExecPass-SingletonDP）：\texttt{ExecutivePass.cpp}}
#include "ExecutivePass.h"

ExecutivePass *ExecutivePass::single_instance = nullptr;   ①

ExecutivePass::~ExecutivePass()
{ 
    single_instance = nullptr;                             ②
}

ExecutivePass *ExecutivePass::obtain(const string holder)  ③
{
    if (single_instance == nullptr)
    {
        single_instance = new ExecutivePass();
    }

    single_instance->holder = holder;
    return single_instance;
}
\end{codeboxt}
\noindent{\fontsize{9bp}{12bp}\selectfont ① 把 single\_instance 初始化为 nullptr，以便能够创建新的单例对象。}\par
\noindent{\fontsize{9bp}{12bp}\selectfont ② 返回指向已存在单例对象的指针。}\par
\noindent{\fontsize{9bp}{12bp}\selectfont ③ 如果尚未存在单例对象，则创建一个并返回指向它的指针。}\par

测试程序表明，我们现在拥有了一个行为正确的单例 \texttt{ExecutivePass} 对象。

\begin{codeboxt}{代码清单 14.5（程序 14.2 ExecPass-SingletonDP）：\texttt{tester.cpp}}
#include <iostream>
#include "ExecutivePass.h"

using namespace std;

ExecutivePass *ExecutivePass::single_instance = nullptr;

int main()
{
    ExecutivePass *pass = ExecutivePass::obtain("Ron");
    cout << "obtain: "; pass->print();

    pass = ExecutivePass::obtain("Sal");
    cout << "obtain: "; pass->print();

    delete pass;
    cout << "delete" << endl;

    pass = ExecutivePass::obtain("Bob");
    cout << "obtain: "; pass->print();

    pass = ExecutivePass::obtain("Flo");
    cout << "obtain: "; pass->print();

    delete pass;
    return 0;
}
\end{codeboxt}

输出如下

\begin{codebox}
obtain: Executive pass now held by Ron, key = 7807    ①
obtain: Executive pass now held by Sal, key = 7807    ②
delete
obtain: Executive pass now held by Bob, key = 6249    ③
obtain: Executive pass now held by Flo, key = 6249    ④
\end{codebox}
\noindent{\fontsize{9bp}{12bp}\selectfont ① Ron 获得通行证。}\par
\noindent{\fontsize{9bp}{12bp}\selectfont ② 把通行证转交给 Sal}\par
\noindent{\fontsize{9bp}{12bp}\selectfont ③ 删除旧通行证之后，Bob 获得一张新通行证。}\par
\noindent{\fontsize{9bp}{12bp}\selectfont ④ 把通行证转交给 Flo}\par

\subsection{14.1.4 单例的通用模型}

图 14.3 展示了单例设计模式的通用模型。回想一下，正是依据设计模式的通用模型，我们才能针对某个架构问题创建定制的解决方案（8.1.3 节）。

\myGraphic{0.739}{images/CH14_F03_Mak.png}{图 14.3 组合设计模式的通用模型。请与图 14.2 对照。私有静态成员变量 \texttt{singleton} 指向 \texttt{Singleton} 对象（如果它存在的话）；否则其值为 \texttt{nullptr}。私有静态成员函数 \texttt{get\_singleton()} 返回指向已存在 \texttt{Singleton} 对象的指针。如果单例对象尚未存在，该函数会创建并返回指向一个新创建 \texttt{Singleton} 对象的指针。}

表 14.1 展示了示例应用如何应用该模式。

\begin{center}
{\bfseries 表 14.1 示例应用所应用的单例方法设计模式}\par\vspace{3bp}
\begin{tabularx}{\textwidth}{XX}
\toprule
设计模式 & 示例应用中的应用 \\
\midrule
类 \texttt{Singleton} & 类 \texttt{ExecutivePass} \\
静态成员变量 \texttt{singleton} & 静态成员变量 \texttt{single\_instance} \\
静态成员函数 \texttt{get\_singleton()} & 静态成员函数 \texttt{obtain()} \\
私有构造函数 \texttt{Singleton()} & 私有构造函数 \texttt{ExecutivePass()} \\
成员变量 \texttt{singleton\_data} & 成员变量 \texttt{key and holder} \\
成员函数 \texttt{singleton\_operation()} & 成员函数 \texttt{print()} \\
\bottomrule
\end{tabularx}
\end{center}

\begin{mySidebar}{单例对象与多线程}

在多线程程序中创建单例对象时必须格外小心。在我们的示例应用中，类 \texttt{ExecutivePass} 的成员函数 \texttt{obtain()} 会先检查单例对象是否已经存在。如果不存在，就创建一个。但是，当单例对象尚不存在时，如果多个线程同时执行这个成员函数，就可能有多个线程都检测到静态成员变量 \texttt{single\_instance} 的值为 \texttt{nullptr} ，于是这些线程各自都会创建一个单例对象。成员变量 \texttt{single\_instance} 将指向最后创建的那个。第 15 章将介绍多线程程序的设计，以及如何避免这类\emph{竞态条件}。
\end{mySidebar}

\section{组合设计模式统一对待单个对象与组合对象}

许多应用把数据保存在树形数据结构中。我们在第 11 章和访问者设计模式中见过树结构。树结构中的数据通常表示“部分—整体”的层级关系。组合设计模式为管理这类层级关系的软件架构提供了模型。

此类层级关系的一个例子，是表示棒球运动员装备用品（provisions）及其每件物品成本的对象（图 14.4）。

\myGraphic{0.600}{images/CH14_F04_Mak.png}{图 14.4 树形数据结构中表示棒球运动员装备用品及其每件物品成本的对象层级。组合对象以灰色阴影表示。}

假设体育部想要一份为棒球运动员置装（outfit）的成本打印输出：

\begin{codebox}
PROVISIONS
    EQUIPMENT
          ball cost: $ 5
           bat cost: $25
         glove cost: $35
    EQUIPMENT total: $65
    UNIFORM
           cap cost: $15
        jersey cost: $25
         pants cost: $35
        FOOTWEAR
             socks cost: $ 5
             shoes cost: $50
        FOOTWEAR total: $55
    UNIFORM total: $130
    sunscreen cost: $ 5
PROVISIONS total: $200

GRAND TOTAL: $200
\end{codebox}

报告中列出了每件物品的成本。它还显示了器材（equipment）、制服（uniform）和鞋类（footwear）这几个组合的小计成本，以及所有物品的总成本。

\subsection{14.2.1 期望的设计特性}

一个用树结构数据来表示“部分—整体”层级关系的应用，应当具备以下特性：

\begin{itemize}
\item \emph{DF 1}——应用应当能够像对待单个部件那样对待部件的组合。
\item \emph{DF 2}——应当能够在运行时添加新的部件以及部件的组合。
\end{itemize}

\subsection{14.2.2 使用组合设计模式之前}

我们成本报告应用的第一个版本不使用组合设计模式。在考察了该版本的设计缺陷之后，我们将使用该模式重构这个应用，并看看该模式如何简化代码。

在我们应用的第一个版本中，类 \texttt{ProvisionItem} 是子类 \texttt{Ball}、\texttt{Bat}、\texttt{Glove}、\texttt{Cap}、\texttt{Jersey}、\texttt{Pants}、\texttt{Socks}、\texttt{Shoes} 和 \texttt{Sunscreen} 的超类。类 \texttt{ProvisionGroup} 是组合类 \texttt{Equipment}、\texttt{Uniform} 和 \texttt{Footwear} 的超类。每个组合类都聚合了 \texttt{ProvisionItem} 子类的一个子集。这些类共同实现了一个树结构的数据层级（图 14.5）。

\myGraphic{0.980}{images/CH14_F05_Mak.png}{图 14.5 类 \texttt{Ball}、\texttt{Bat}、\texttt{Glove}、\texttt{Cap}、\texttt{Jersey}、\texttt{Pants}、\texttt{Socks}、\texttt{Shoes} 和 \texttt{Sunscreen} 是类 \texttt{ProvisionItem} 的子类，每个子类都实现了 \texttt{get\_cost()} 成员函数。组合类 \texttt{Equipment}、\texttt{Uniform} 和 \texttt{Footwear} 是 \texttt{ProvisionGroup} 的子类，每个都聚合了 \texttt{ProvisionItem} 子类的一个子集。}

类 \texttt{ProvisionItem} 是子类 \texttt{Ball}、\texttt{Bat}、\texttt{Glove}、\texttt{Cap}、\texttt{Jersey}、\texttt{Pants}、\texttt{Socks}、\texttt{Shoes} 和 \texttt{Sunscreen} 的超类。该超类实现了公共的 \texttt{get\_id()}、\texttt{get\_cost()} 和 \texttt{print()} 成员函数，并拥有私有成员变量 \texttt{id} 和 \texttt{cost}。

\begin{codeboxt}{代码清单 14.6（程序 14.3 CostReport）：\texttt{ProvisionItem.h}（使用设计模式之前）}
#include <iostream>
#include <iomanip>
#include <string>

using namespace std;

class ProvisionItem
{
public:
    ProvisionItem(const string id, const double c)
        : id(id), cost(c) {}

    string get_id()   const { return id; }
    double get_cost() const { return cost; }

    void print()
    {
        cout << setw(6) << get_id() << " cost: $"
             << setw(2) << get_cost() << endl;
    }

private:
    string id;
    double cost;
};
\end{codeboxt}

\texttt{ProvisionItem} 的子类 \texttt{Ball}、\texttt{Bat} 和 \texttt{Glove} 表示器材物品。每个子类的构造函数都会调用超类 \texttt{ProvisionItem} 的构造函数，并传入自己的 \texttt{id} 和 \texttt{cost}。

\begin{codeboxt}{代码清单 14.7（程序 14.3 CostReport）：\texttt{Equipment.h}（使用设计模式之前）}
#include "ProvisionItem.h"

using namespace std;

class Ball : public ProvisionItem
{
public:
    Ball(const double cost) : ProvisionItem("ball", cost ) {}
};

class Bat : public ProvisionItem
{
public:
    Bat(const double cost) : ProvisionItem("bat", cost ) {}
};

class Glove : public ProvisionItem
{
public:
    Glove(const double cost) : ProvisionItem("glove", cost ) {}
};
\end{codeboxt}

子类 \texttt{Cap}、\texttt{Jersey} 和 \texttt{Pants} 各自表示一件制服物品。

\begin{codeboxt}{代码清单 14.8（程序 14.3 CostReport）：\texttt{Uniform.h}（使用设计模式之前）}
#include "ProvisionItem.h"

using namespace std;

class Cap : public ProvisionItem
{
public:
    Cap(const double cost) : ProvisionItem("cap", cost ) {}
};

class Jersey : public ProvisionItem
{
public:
    Jersey(const double cost) : ProvisionItem("jersey", cost ) {}
};

class Pants : public ProvisionItem
{
public:
    Pants(const double cost) : ProvisionItem("pants", cost ) {}
};
\end{codeboxt}

子类 \texttt{Socks} 和 \texttt{Shoes} 各自表示一件鞋类物品。

\begin{codeboxt}{代码清单 14.9（程序 14.3 CostReport）：\texttt{Footwear.h}（使用设计模式之前）}
#include "ProvisionItem.h"

using namespace std;

class Socks : public ProvisionItem
{
public:
    Socks(const double cost) : ProvisionItem("socks", cost) {}
};

class Shoes : public ProvisionItem
{
public:
    Shoes(const double cost) : ProvisionItem("shoes", cost ) {}
};
\end{codeboxt}

而子类 Sunscreen 表示一件防晒霜物品。

\begin{codeboxt}{代码清单 14.10（程序 14.3 CostReport）：\texttt{Sunscreen.h}（使用设计模式之前）}
#include "ProvisionItem.h"

using namespace std;

class Sunscreen : public ProvisionItem
{
public:
    Sunscreen(const double cost) 
        : ProvisionItem("sunscreen", cost) {}
};
\end{codeboxt}

类 \texttt{Equipment}、\texttt{Uniform} 和 \texttt{Footwear} 是超类 \texttt{ProvisionGroup} 的子类。该超类提供了公共成员变量 \texttt{provisions}，它是一个存放 \texttt{ProvisionItem} 对象的 vector。\texttt{add()} 成员函数把 \texttt{ProvisionItem} 对象追加到该 vector 中，\texttt{remove()} 成员函数则从该 vector 中移除一个 \texttt{ProvisionItem} 对象。

\begin{codeboxt}{代码清单 14.11（程序 14.3 CostReport）：\texttt{ProvisionGroup.h}（使用设计模式之前）}
#include <vector>
#include <algorithm>
#include "ProvisionItem.h"

using namespace std;

class ProvisionGroup
{
public:
    ProvisionGroup(string id) : id(id) {}
    ~ProvisionGroup()
    {
        for (ProvisionItem *item : provisions) { delete pi; };
    }

    string get_id() const { return id; }

    vector<ProvisionItem *> get_provisions() const
    {
        return provisions;
    }

    void add(ProvisionItem *item) { provisions.push_back(item); }

    void remove(ProvisionItem *item)
    {
        auto pos = find(provisions.begin(), provisions.end(), item);
        if (pos != provisions.end()) provisions.erase(pos);
    }

private:
    vector<ProvisionItem *> provisions;
};

class Equipment : public ProvisionGroup    ①
{
public:
    Equipment() : ProvisionGroup("EQUIPMENT") {}
};

class Uniform : public ProvisionGroup      ①
{
public:
    Uniform() : ProvisionGroup("UNIFORM") {}
};

class Footwear : public ProvisionGroup     ①
{
public:
    Footwear() : ProvisionGroup("FOOTWEAR") {}
};
\end{codeboxt}
\noindent{\fontsize{9bp}{12bp}\selectfont ① 超类 ProvisionGroup 的子类}\par

类 \texttt{CostReport} 创建数据并生成报告。

\begin{codeboxt}{代码清单 14.12（程序 14.3 CostReport）：\texttt{CostReport.h}（使用设计模式之前）}
#include "ProvisionGroups.h"
#include "Sunscreen.h"

class CostReport
{
public:
    CostReport();
    ~CostReport();

    void generate_report();

private:
    Equipment *equipment;
    Uniform   *uniform;
    Footwear  *footwear;
    Sunscreen *sunscreen;
};
\end{codeboxt}

构造函数构建数据层级。它创建各个 \texttt{ProvisionItem} 对象，并把它们添加到相应的 \texttt{Equipment}、\texttt{Uniform} 和 \texttt{Footwear} 对象中。它还创建独立的 \texttt{Sunscreen} 对象。

\begin{codeboxt}{代码清单 14.13（程序 14.3 CostReport）：\texttt{CostReport.cpp}（使用设计模式之前）}
#include <iostream>
#include <iomanip>
#include "ProvisionItem.h"
#include "ProvisionGroups.h"
#include "Equipment.h"
#include "Uniform.h"
#include "Footwear.h"
#include "Sunscreen.h"
#include "CostReport.h"

using namespace std;

CostReport::CostReport()
    : equipment(new Equipment()), uniform(new Uniform()),
      footwear(new Footwear()), sunscreen(new Sunscreen(5))
{
    equipment->add(new Ball(5));
    equipment->add(new Bat(25));
    equipment->add(new Glove(35));

    uniform->add(new Cap(15));
    uniform->add(new Jersey(25));
    uniform->add(new Pants(35));

    footwear->add(new Socks(5));
    footwear->add(new Shoes(50));
}

CostReport::~CostReport()
{
    delete equipment;
    delete uniform;
    delete footwear;
    delete sunscreen;
}

...
\end{codeboxt}

成员函数 \texttt{generate\_report()} 生成成本报告。它先计算器材、制服和鞋类物品的总成本，然后计算总体的总成本。

\begin{codeboxt}{代码清单 14.14（程序 14.3 CostReport）：\texttt{CostReport.cpp}（使用设计模式之前）}
...

void CostReport::generate_report()
{
    double equipment_cost = 0;                            ①
    for (ProvisionItem *pi : equipment->get_provisions()) ①
    {                                                     ①
        equipment_cost += pi->get_cost();                 ①
    }                                                     ①
                                                          ①
    double uniform_cost = 0;                              ①
    for (ProvisionItem *pi : uniform->get_provisions())   ①
    {                                                     ①
        uniform_cost += pi->get_cost();                   ①
    }                                                     ①
                                                          ①
    double footwear_cost = 0;                             ①
    for (ProvisionItem *pi : footwear->get_provisions())  ①
    {                                                     ①
        Footwear_cost += pi->get_cost();                  ①
    }                                                     ①
                                                          ①
    uniform_cost += footwear_cost;                        ①

    cout << "PROVISIONS" << endl;

    cout << "    " << equipment->get_id() << endl;
    for (ProvisionItem *pi : equipment->get_provisions()) ②
    {                                                     ②
        cout << "        ";                               ②
        pi->print();                                      ②
    }                                                     ②
    cout << "    " << equipment->get_id() << " total: $"
         << setw(2) << equipment_cost << endl;

    cout << "    " << uniform->get_id() << endl;
    for (ProvisionItem *pi : uniform->get_provisions())   ②
    {                                                     ②
        cout << "        ";                               ②
        pi->print();                                      ②
    }                                                     ②

    cout << "        " << footwear->get_id() << endl;
    for (ProvisionItem *pi : footwear->get_provisions())  ②
    {                                                     ②
        cout << "            ";                           ②
        pi->print();                                      ②
    }                                                     ②
    cout << "        " << footwear->get_id() << " total: $"
         << setw(2) << footwear_cost << endl;

    cout << "    " << uniform->get_id() << " total: $"
         << setw(2) << uniform_cost << endl;

    cout << "    ";
    sunscreen->print();

    double provisions_total = equipment_cost             ③
                            + uniform_cost               ③
                            + sunscreen->get_cost();     ③
    cout << "PROVISIONS total: $" << provisions_total << endl;

    cout << endl;
    cout << "GRAND TOTAL: $" << provisions_total << endl;
}
\end{codeboxt}
\noindent{\fontsize{9bp}{12bp}\selectfont ① 计算每个组合 ProvisionGroup 的成本。}\par
\noindent{\fontsize{9bp}{12bp}\selectfont ② 打印各个 ProvisionItem 对象的成本。}\par
\noindent{\fontsize{9bp}{12bp}\selectfont ③ 计算总体成本。}\par

测试程序生成了期望的成本报告。

\begin{codeboxt}{代码清单 14.15（程序 14.3 CostReport）：\texttt{tester.cpp}}
#include "CostReport.h"

int main()
{
    CostReport report;
    report.generate_report();

    return 0;
}
\end{codeboxt}

\myGraphic{0.854}{images/CH14_F05_UN03_Mak.png}{}

确实，这个版本的应用存在以下问题

\begin{itemize}
\item \emph{报告生成器复杂}——类 \texttt{CostReport} 的成员函数 \texttt{generate\_report()} 很复杂。如果改变数据的组织方式，它将难以修改。
\item \emph{单个对象与组合对象被区别对待}——我们希望树中所有层级的对象，无论是单个对象还是组合对象，都被同等对待。而成员函数 \texttt{generate\_report()} 却区别对待它们。
\item \emph{报告生成器被硬编码}——成员函数 \texttt{generate\_report()} 被硬编码为只能打印那一份报告。
\item \emph{树不完整}——我们没有构建完整的树形数据结构。例如，鞋类没有成为制服的一部分。对于本应用生成成本报告而言，不完整的树不是问题，但其他应用可能需要完整的树。
\end{itemize}

\subsection{14.2.3 使用组合设计模式之后}

我们可以把成本报告应用的第一个版本重构为仿照组合设计模式的架构。组合设计模式的目标，是把树形数据结构中的单个对象与对象组合同等对待。我们希望数据树的所有叶类 \texttt{Ball}、\texttt{Bat}、\texttt{Glove}、\texttt{Cap}、\texttt{Jersey}、\texttt{Pants}、\texttt{Socks}、\texttt{Shoes} 和 \texttt{Sunscreen}，以及组合类 \texttt{Equipment}、\texttt{Uniform} 和 \texttt{Footwear}，都呈现相同的接口。

\begin{mySidebar}{组合设计模式}

“将对象组合成树形结构以表示‘部分—整体’层级。组合让客户端统一地对待单个对象和对象组合”（GoF 163）。
\end{mySidebar}

为了使用这个模型，我们还把类 \texttt{ProvisionItem} 设为类 \texttt{ProvisionGroup} 的超类。这样一来，组合类 \texttt{Equipment}、\texttt{Uniform} 和 \texttt{Footwear} 最终就成了 \texttt{ProvisionItem} 的子类。类 \texttt{Ball}、\texttt{Bat}、\texttt{Glove}、\texttt{Cap}、\texttt{Jersey}、\texttt{Pants}、\texttt{Socks}、\texttt{Shoes} 和 \texttt{Sunscreen} 仍为 \texttt{ProvisionItem} 的子类，而类 \texttt{ProvisionGroup} 维护一个存放 \texttt{ProvisionItem} 对象的 vector（图 14.6）。

\myGraphic{0.980}{images/CH14_F06_Mak.png}{图 14.6 在由组合设计模式建模的成本报告应用第二个版本中，类 \texttt{ProvisionGroup} 本身就是类 \texttt{ProvisionItem} 的子类。这使得类 \texttt{Equipment}、\texttt{Uniform} 和 \texttt{Footwear} 最终成为 \texttt{ProvisionItem} 的子类。图中灰色部分相对于图 14.5 在逻辑上没有变化。}

类 \texttt{ProvisionItem} 和 \texttt{ProvisionGroup} 将拥有相同的成员函数。因此，类 \texttt{ProvisionItem} 增加了成员函数 \texttt{get\_provisions()}、\texttt{add()} 和 \texttt{remove()}。类 \texttt{ProvisionGroup} 继承成员函数 \texttt{get\_cost()} 和 \texttt{print()}。类 \texttt{ProvisionItem} 为其所有成员函数提供了默认实现（代码清单 14.16）。超类 \texttt{ProvisionItem} 的成员函数 \texttt{get\_id()} 和 \texttt{get\_cost()} 没有改变——我们希望子类 \texttt{Ball}、\texttt{Bat}、\texttt{Glove}、\texttt{Cap}、\texttt{Jersey}、\texttt{Pants}、\texttt{Socks}、\texttt{Shoes} 和 \texttt{Sunscreen} 继续继承并使用这些函数。成员函数 \texttt{print()} 将接受一个整数参数，用于决定在打印行开头缩进多少个空格（如果有的话）。

由于成员函数 \texttt{get\_provisions()}、\texttt{add()} 和 \texttt{remove()} 对单个对象没有意义，因此每个函数的默认实现都会抛出 \texttt{InvalidOperation} 异常。组合类 \texttt{Equipment}、\texttt{Uniform} 和 \texttt{Footwear} 会重写这些成员函数以实现有用的行为，因此它们是虚函数。

\begin{codeboxt}{代码清单 14.16（程序 14.4 CostReport-CompositeDP）：\texttt{ProvisionItem.h}}
#include <iostream>
#include <iomanip>
#include <string>
#include <vector>

using namespace std;

class InvalidOperation {};

class ProvisionItem
{
public:
    ProvisionItem(const string id) : id(id), cost(0) {}

    ProvisionItem(const string id, const double c)
        : id(id), cost(c) {}

    virtual ~ProvisionItem() {}

    string get_id() const { return id; }

    virtual double get_cost() const { return cost; }

    virtual void print(const int indentation)
    {
        for (int i = 0; i < indentation; i++) cout << "    "; ①

        cout << setw(6) << get_id() << " cost: $"
             << setw(2) << get_cost() << endl;
    }

    virtual vector<ProvisionItem *> get_provisions() const
    {
        throw new InvalidOperation();                         ②
    }

    virtual void add(ProvisionItem *item)
    {
        throw new InvalidOperation();                         ②
    }

    virtual void remove(ProvisionItem *item)
    {
        throw new InvalidOperation();                         ②
    }

private:
    string id;
    double cost;
};
\end{codeboxt}
\noindent{\fontsize{9bp}{12bp}\selectfont ① 打印前缩进。}\par
\noindent{\fontsize{9bp}{12bp}\selectfont ② 单个对象的默认行为}\par

\myGraphic{0.833}{images/CH14_F06_UN04_Mak.png}{}

一致性使我们仿照组合设计模式建模的应用能够像对待诸如 \texttt{Footwear} 对象这样的组合对象那样，对待诸如 \texttt{Ball} 对象这样的单个对象。报告生成代码将更为简单——它无需知道正在处理的是单个对象还是组合对象。

类 \texttt{ProvisionGroup} 的公有构造函数和受保护构造函数都会调用其超类 \texttt{ProvisionItem} 的构造函数。如以下代码清单所示，这个组合类重写了成员函数 \texttt{get\_cost()}、\texttt{print()}、\texttt{get\_provisions()}、\texttt{get\_provisions()}、\texttt{add()} 和 \texttt{remove()}，以实现适合组合对象的行为。成员函数 \texttt{get\_cost()} 和 \texttt{print()} 都会遍历 \texttt{provisions} vector 中的 \texttt{ProvisionItem} 对象，分别对每个对象调用 \texttt{get\_cost()} 和 \texttt{print()}。

\begin{codeboxt}{代码清单 14.17（程序 14.4 CostReport-CompositeDP）：\texttt{ProvisionGroup.h}（第 1 部分，共 2 部分）}
#include <iostream>
#include <iomanip>
#include <string>
#include <vector>
#include <algorithm>
#include "ProvisionItem.h"

using namespace std;

class ProvisionGroup : public ProvisionItem
{
public:
    ProvisionGroup() : ProvisionItem("PROVISIONS", 0) {}

    ~ProvisionGroup()
    {
        for (ProvisionItem *item : provisions) { delete item; };
    }

    double get_cost() const override
    {
        double cost = 0;                                    ①
        for (ProvisionItem *item : provisions)              ①
        {                                                   ①
            cost += item->get_cost();                       ①
        }                                                   ①

        return cost;
    }

    void print(const int indentation) override
    {
        for (int i = 0; i < indentation; i++) cout << "    ";
        cout << get_id() << endl;

        for (ProvisionItem *item : provisions)              ②
        {                                                   ②
            item->print(indentation + 1);                   ②
        }                                                   ②

        for (int i = 0; i < indentation; i++) cout << "    ";
        cout << get_id() << " total: $" << get_cost() << endl;
    }

    vector<ProvisionItem *> get_provisions() const override ③
    {
        return provisions;
    }

    void add(ProvisionItem *item) override                  ③
    {
        provisions.push_back(item);
    }

    void remove(ProvisionItem *item) override               ③
    {
        auto pos = find(provisions.begin(), provisions.end(), item);
        if (pos != provisions.end()) provisions.erase(pos);
    }

protected:
    ProvisionGroup(const string id) : ProvisionItem(id) {}

private:
    vector<ProvisionItem *> provisions;
};
...
\end{codeboxt}
\noindent{\fontsize{9bp}{12bp}\selectfont ① 遍历 ProvisionItem 对象以计算这些物品的总成本。}\par
\noindent{\fontsize{9bp}{12bp}\selectfont ② 遍历 ProvisionItem 对象以打印每个对象。}\par
\noindent{\fontsize{9bp}{12bp}\selectfont ③ 为 ProvisionGroup 对象重写的成员函数}\par

组合子类 \texttt{Equipment}、\texttt{Uniform} 和 \texttt{Footwear} 的每个构造函数都会调用超类 \texttt{ProvisionGroup} 的受保护构造函数。

\begin{codeboxt}{代码清单 14.18（程序 14.4 CostReport-CompositeDP）：\texttt{ProvisionGroup.h}（第 2 部分，共 2 部分）}
...

class Equipment : public ProvisionGroup
{
public:
    Equipment() : ProvisionGroup("EQUIPMENT") {}
};

class Uniform : public ProvisionGroup
{
public:
    Uniform() : ProvisionGroup("UNIFORM") {}
};

class Footwear : public ProvisionGroup
{
public:
    Footwear() : ProvisionGroup("FOOTWEAR") {}
};
\end{codeboxt}

类 \texttt{CostReport} 增加了一个成员变量 \texttt{provisions}。

\begin{codeboxt}{代码清单 14.19（程序 14.4 CostReport-CompositeDP）：\texttt{CostReport.h}（第 2 部分，共 2 部分）}
#include "ProvisionGroup.h"
#include "Sunscreen.h"

class CostReport
{
public:
    CostReport();
    ~CostReport();

    void generate_report();

private:
    ProvisionGroup provisions;

    Equipment *equipment;
    Uniform   *uniform;
    Footwear  *footwear;
    Sunscreen *sunscreen;
};
\end{codeboxt}

类 \texttt{CostReport} 的构造函数构建了图 14.4 的整个数据树。类 \texttt{Footwear} 成为类 \texttt{Uniform} 下的一个子集合。

\begin{codeboxt}{代码清单 14.20（程序 14.4 CostReport-CompositeDP）：\texttt{CostReport.cpp}}
#include <iostream>
#include <iomanip>
#include "ProvisionItem.h"
#include "ProvisionGroup.h"
#include "Equipment.h"
#include "Uniform.h"
#include "Footwear.h"
#include "Sunscreen.h"
#include "CostReport.h"

using namespace std;

CostReport::CostReport()
    : equipment(new Equipment()), uniform(new Uniform()),
      footwear(new Footwear()), sunscreen(new Sunscreen)
{
    equipment->add(new Ball(5));
    equipment->add(new Bat(25));
    equipment->add(new Glove(35));

    uniform->add(new Cap(15));
    uniform->add(new Jersey(25));
    uniform->add(new Pants(35));

    footwear->add(new Socks(5));
    footwear->add(new Shoes(50));

    uniform->add(footwear);     ①

    provisions.add(equipment);
    provisions.add(uniform);
    provisions.add(sunscreen);
}

CostReport::~CostReport()
{
    delete equipment;
    delete uniform;
    delete footwear;
    delete sunscreen;
}
void CostReport::generate_report()
{
    provisions.print(0);

    cout << endl;
    cout << "GRAND TOTAL: $" << setw(3) << provisions.get_cost()
         << endl;
}
\end{codeboxt}
\noindent{\fontsize{9bp}{12bp}\selectfont ① 让 Footwear 成为 Uniform 的子集合。}\par

\texttt{ProvisionItem} 的子类 \texttt{Ball}、\texttt{Bat}、\texttt{Glove}、\texttt{Cap}、\texttt{Jersey}、\texttt{Pants}、\texttt{Socks}、\texttt{Shoes} 和 \texttt{Sunscreen} 在这个版本的应用中没有变化，测试程序也没有变化。输出是相同的。

我们使用组合设计模式所获得的主要好处有

\begin{itemize}
\item \emph{一致性}——单个对象的类与组合对象的类共享同一个接口。因此，我们的代码能够像对待组合对象那样对待单个对象。
\item \emph{简单性}——对树中各组件执行操作（如计算成本和打印）的代码更简单。
\item \emph{灵活性}——添加新的 \texttt{ProvisionItem} 和 \texttt{ProvisionGroup} 对象或移除已有对象都不难。
\end{itemize}

\subsection{14.2.4 组合的通用模型}

图 14.7 展示了组合设计模式的通用模型。回想一下，正是依据设计模式的通用模型，我们才能针对某个架构问题创建定制的解决方案（8.1.3 节）……

\myGraphic{0.980}{images/CH14_F07_Mak.png}{图 14.7 组合设计模式的通用模型。请与图 14.6 对照。单个对象类 \texttt{Leaf} 和 \texttt{Composite} 类都继承自超类 \texttt{Component}。}

表 14.2 展示了示例应用如何应用该模式。

\begin{center}
{\bfseries 表 14.2 示例应用所应用的组合设计模式}\par\vspace{3bp}
\begin{tabularx}{\textwidth}{XX}
\toprule
设计模式 & 示例应用中的应用 \\
\midrule
超类 \texttt{Component} & 超类 \texttt{ProvisionItem} \\
类 \texttt{Composite} & 类 \texttt{ProvisionGroup} \\
类 \texttt{Leaf} & 类 \texttt{Ball}、\texttt{Bat}、\texttt{Glove}、\texttt{Cap}、\texttt{Jersey}、\texttt{Pants}、\texttt{Socks}、\texttt{Shoes} 和 \texttt{Sunscreen} \\
\bottomrule
\end{tabularx}
\end{center}

\myGraphic{0.980}{images/CH14_F07_UN05_Mak.png}{}

\section{装饰器设计模式动态地添加对象职责}

类有成员变量，其运行时值就是对象的各种属性（attribute）。正如第 2 章所强调的，管理对象属性是应用设计的一个重要方面。类也有成员函数，用于控制其对象的运行时行为。成员变量和成员函数实现了类的职责，良好的设计会为每个类分配一个主要职责。但是，如果我们需要灵活地为类的对象动态添加更多职责，该怎么办呢？装饰器设计模式为灵活管理对象职责的应用架构提供了模型。

作为此类应用的一个例子，假设除了通过售票机售出的基础门票之外，体育部还提供特种增强票（enhanced ticket）的在线销售。门票增强项可以包括赛前派对邀请、VIP 座位和饮品券。我们可以把这些增强项实现为门票的属性，而每个增强项都需要一种行为：把它自身的成本加到门票的基础价格上。我们可以把每个门票增强项都视为门票的一项额外职责。

\subsection{14.3.1 期望的设计特性}

一个应用若要能够为已有的类动态添加额外职责，应当具备以下特性：

\begin{itemize}
\item \emph{DF 1}——我们应当无需修改要被添加新职责的那个类。
\item \emph{DF 2}——我们应当能够在运行时以任意顺序和任意组合动态添加任意数量的职责。
\item \emph{DF 3}——该类不应依赖这些职责的具体实现。
\item \emph{DF 4}——应当能够在不修改拥有这些职责的类的情况下，添加、删除或修改职责。
\item \emph{DF 5}——使用某个能够获得额外职责的对象的代码，不应根据额外职责的数量（包括没有）、顺序或组合而对该对象有任何不同的对待。
\end{itemize}

\subsection{14.3.2 使用装饰器设计模式之前}

我们将开发增强门票应用的两个版本。第一个版本以一种直白的方式实现门票的属性和行为，即把它们作为 \texttt{Ticket} 类的成员变量和成员函数。但那样实现会有一些重大缺陷。随后，我们将仿照装饰器设计模式来为应用的第二个版本设计建模，这将大幅提高管理额外门票行为的灵活性。

我们门票应用的第一个版本只是把每个门票增强项及其成本作为 \texttt{Ticket} 类的一个属性（图 14.8）。

\myGraphic{0.375}{images/CH14_F08_Mak.png}{图 14.8 我们把对基础门票的增强项实现为 \texttt{Ticket} 类的成员变量。每个增强项都有一个成本。}

类 \texttt{Ticket} 的成员变量实现了门票的增强项及其成本。如以下代码清单所示，成员函数 \texttt{get\_cost()} 计算并返回一张门票的总成本。\texttt{Ticket} 对象的职责是管理门票的增强项并计算其总成本。

\begin{codeboxt}{代码清单 14.21（程序 14.5 Enhanced）：\texttt{Ticket.h}（使用设计模式之前）}
#include <iostream>
#include <iomanip>

using namespace std;

class Ticket
{
public:
    Ticket(const bool party, const bool vip, const int coupons)
       : pregame_party(party), vip_seating(vip),
         drink_coupons(coupons) {}

    double get_cost() const
    {
       double total = BASE_PRICE;

       if (pregame_party) total += PARTY_PRICE;
       if (vip_seating)   total += VIP_PRICE;

       return total + drink_coupons*COUPON_PRICE;
    }
    friend ostream& operator <<(ostream& ostr, const Ticket& ticket);

private:
    bool pregame_party;                ①
    bool vip_seating;                  ①
    int  drink_coupons;                ①

    double const BASE_PRICE   = 30;    ②
    double const PARTY_PRICE  = 25;    ②
    double const VIP_PRICE    = 20;    ③
    double const COUPON_PRICE =  5;    ③
};

inline ostream& operator <<(ostream& ostr, const Ticket& ticket)
{
    ostr << setw(24) << "base ticket price: $"
         << setw(2)  << ticket.BASE_PRICE << endl;

    if (ticket.pregame_party)
    {
        ostr << setw(24) << "pregame party price: $"
             << setw(2) << ticket.PARTY_PRICE << endl;
    }

    if (ticket.vip_seating)
    {
        ostr << setw(24) << "VIP seating price: $"
             << setw(2)  << ticket.VIP_PRICE << endl;
    }

    int count = ticket.drink_coupons;

    if (count > 0)
    {
        ostr << setw(24) << "drink coupon price: $"
             << setw(2)  << count*ticket.COUPON_PRICE
             << " = " << count << " x $"
             << ticket.COUPON_PRICE << endl;
    }

    return ostr;
}
\end{codeboxt}
\noindent{\fontsize{9bp}{12bp}\selectfont ① 增强项}\par
\noindent{\fontsize{9bp}{12bp}\selectfont ② 成本}\par
\noindent{\fontsize{9bp}{12bp}\selectfont ③ 成本}\par

一个测试程序创建并打印一些带有增强项的示例门票。

\begin{codeboxt}{代码清单 14.22（程序 14.5 Enhanced）：\texttt{tester.cpp}（使用设计模式之前）}
#include <iostream>
#include <string>
#include "Ticket.h"

using namespace std;

void print(const string& name, const Ticket& ticket);

int main()
{
    Ticket ticket1(true, true, 2);
    print("John", ticket1);

    Ticket ticket2(true, false, 3);
    print("Mary", ticket2);

    Ticket ticket3(false, false, 0);
    print("Leslie", ticket3);

    Ticket ticket4(false, true, 1);
    print("Sidney", ticket4);

    return 0;
}

void print(const string& name, const Ticket& ticket)
{
    cout << endl;
    cout << name << "'s ticket:" << endl;
    cout << ticket;
    cout << "TOTAL COST: $" << ticket.get_cost() << endl;
}
\end{codeboxt}

输出如下

\begin{codebox}
John's ticket:
    base ticket price: $30
  pregame party price: $25
    VIP seating price: $20
   drink coupon price: $10 = 2 x $5
TOTAL COST: $85

Mary's ticket:
    base ticket price: $30
  pregame party price: $25
   drink coupon price: $15 = 3 x $5
TOTAL COST: $70

Leslie's ticket:
    base ticket price: $30
TOTAL COST: $30

Sidney's ticket:
    base ticket price: $30
    VIP seating price: $20
   drink coupon price: $ 5 = 1 x $5
TOTAL COST: $55
\end{codebox}

\myGraphic{0.679}{images/CH14_F08_UN06_Mak.png}{}

确实，缺乏灵活性是这种设计的主要缺陷：

\begin{itemize}
\item \emph{代码不灵活}——由于把门票增强项的种类硬编码了，要想添加、移除或修改增强项，就必须修改 \texttt{Ticket} 类。
\item \emph{增强项的实现未被隐藏}——类 \texttt{Ticket} 对门票增强项的实现方式了如指掌，并且会根据增强项的不同而区别对待门票（即计算它们的成本）。
\end{itemize}

\subsection{14.3.3 使用装饰器设计模式之后}

装饰器设计模式把每个门票增强项都视为一个装饰，它在概念上包裹了基础门票以及此前的任何增强项（图 14.9）。

\myGraphic{0.980}{images/CH14_F09_Mak.png}{图 14.9 一张增强门票在概念上被任意数量的赛前派对、VIP 座位和饮品券装饰器所包裹。每个装饰器都包裹基础门票以及此前的任何包裹器。为了计算一张增强门票的成本，每个装饰器都会调用它所包裹对象（另一个装饰器或基础门票）的 \texttt{get\_cost()} 成员函数。每个返回值都是该装饰器的成本加上它所包裹对象的总成本。}

图中展示了一张用赛前派对增强项装饰的基础门票，它们又被一个 VIP 座位增强项和两张饮品券进一步装饰。每个装饰器都为门票添加了这样一项职责：把装饰器的成本计入门票总价。每个装饰器和门票都有一个 \texttt{get\_cost()} 函数。每个装饰器的函数都返回该装饰器的成本加上它所包裹对象的总成本。

为了实现这种嵌套包裹的架构，我们把每个装饰器都做成对象，并按它们相互包裹的顺序把这些装饰器对象链接起来。链中的最后一个装饰器与基础门票对象相连（图 14.10）。因此，门票本身被包裹得最深。

\myGraphic{0.641}{images/CH14_F10_Mak.png}{图 14.10 链中除最后一个之外的每个装饰器对象都指向它所包裹的装饰器对象。链中的最后一个装饰器对象指向门票对象。}

还有一项额外要求：无论一张门票因其装饰而拥有多少额外职责，它仍然是一张门票。我们的应用代码应当像对待未装饰的基础门票那样对待已装饰的门票，不作任何区分。为满足这一要求，我们需要让基础门票类和各个装饰器类都成为 \texttt{Ticket} 超类的子类（图 14.11）。我们还可以设计一个 \texttt{Decorator} 超类，其子类就是各个装饰器类。它的成员变量将包含指向下一个 \texttt{Decorator} 包裹器对象或 \texttt{Ticket} 对象的链接。\texttt{Decorator} 本身将是 \texttt{Ticket} 的子类。

\begin{mySidebar}{装饰器设计模式}

“动态地给一个对象附加额外的职责。就扩展功能而言，装饰器提供了一种比继承更灵活的替代方案”（GoF 175）。
\end{mySidebar}

\myGraphic{0.980}{images/CH14_F11_Mak.png}{图 14.11 抽象超类 \texttt{Ticket} 的子类表示基础门票或被一或多个装饰器包裹的门票。\texttt{Decorator} 超类有一个 \texttt{ticket} 链接，指向下一个 \texttt{Decorator} 对象或 \texttt{BaseTicket} 对象。\texttt{Party}、\texttt{VIP} 和 \texttt{Coupon} 是装饰器。它们中一或多个对象可以包裹一个 \texttt{BaseTicket} 对象。每个 \texttt{Decorator} 对象都添加了这样一项职责：把自身的成本计入门票总成本。}

抽象类 \texttt{Ticket} 是类 \texttt{BaseTicket} 和超类 \texttt{Decorator} 的超类。

\begin{codeboxt}{代码清单 14.23（程序 14.6 Enhanced-DecoratorDP）：\texttt{Ticket.h}}
#include <iostream>
#include <iomanip>
#include <string>

using namespace std;

class Ticket
{
public:
    Ticket(const string desc, const double pr)
        : description(desc), price(pr)
    {
        cout << setw(17) << description                     ①
             << " price: $" << setw(2) << price << endl;    ①
    }

    virtual ~Ticket() {}

    virtual double get_cost() const = 0;

    string get_description() const { return description; }

protected:
    string description;
    double price;
};
\end{codeboxt}
\noindent{\fontsize{9bp}{12bp}\selectfont ① 创建 Ticket 对象时打印它。}\par

类 \texttt{BaseTicket} 的每个对象都表示尚未被任何装饰的门票。它的成本就是 \texttt{BASE\_PRICE}。

\begin{codeboxt}{代码清单 14.24（程序 14.6 Enhanced-DecoratorDP）：\texttt{BaseTicket.h}}
#include "Ticket.h"

const double BASE_PRICE = 30;

class BaseTicket : public Ticket
{
public:
    BaseTicket() : Ticket("base ticket", BASE_PRICE) {}

    double get_cost() const override { return BASE_PRICE; }
};
\end{codeboxt}

超类 \texttt{Decorator} 表示门票的装饰。每个装饰器对象都添加了这样一项职责：把自身的成本计入门票总成本。成员函数 \texttt{get\_cost()} 返回它自身的价格加上它所包裹对象的成本。成员变量 \texttt{ticket} 指向链中的下一个 \texttt{Ticket} 对象，无论是 \texttt{BaseTicket} 对象还是另一个 \texttt{Decorator} 对象。

\begin{codeboxt}{代码清单 14.25（程序 14.6 Enhanced-DecoratorDP）：\texttt{Decorator.h}}
#include <string>
#include "Ticket.h"

using namespace std;

class Decorator : public Ticket
{
public:
    Decorator(const string description, const double price,
              Ticket *tckt)
        : Ticket(description, price), ticket(tckt) {}

    virtual ~Decorator() { delete ticket; }

    double get_cost() const override
    {
        return price + ticket->get_cost(); ①
    }

private:
    Ticket *ticket;                         ②
};
\end{codeboxt}
\noindent{\fontsize{9bp}{12bp}\selectfont ① 返回这个 Ticket 对象的价格加上它所包裹对象的成本。}\par
\noindent{\fontsize{9bp}{12bp}\selectfont ② 这个对象所包裹的 Ticket 对象}\par

类 \texttt{Party} 是门票的赛前派对装饰。

\begin{codeboxt}{代码清单 14.26（程序 14.6 Enhanced-DecoratorDP）：\texttt{Party.h}}
#include "Decorator.h"

const double PARTY_PRICE = 25;

class Party : public Decorator
{
public:
    Party(Ticket *ticket)
        : Decorator("pregame party", PARTY_PRICE, ticket) {}
};
\end{codeboxt}

类 \texttt{VIP} 是门票的 VIP 座位装饰。

\begin{codeboxt}{代码清单 14.27（程序 14.6 Enhanced-DecoratorDP）：\texttt{VIP.h}}
#include "Decorator.h"

const double VIP_PRICE = 20;

class VIP : public Decorator
{
public:
    VIP(Ticket *ticket)
        : Decorator("VIP seating", VIP_PRICE, ticket) {}
};
\end{codeboxt}

最后，类 \texttt{Coupon} 是门票的饮品券装饰。

\begin{codeboxt}{代码清单 14.28（程序 14.6 Enhanced-DecoratorDP）：\texttt{Coupon.h}}
#include "Decorator.h"
const double COUPON_PRICE = 5;
class Coupon : public Decorator
{
public:
    Coupon(Ticket *ticket)
        : Decorator("drink coupon", COUPON_PRICE, ticket) {}
};
\end{codeboxt}

测试程序创建 \texttt{BaseTicket} 对象，并通过以任意顺序、任意数量和组合的 \texttt{Decorator} 对象包裹每张门票来装饰它。如果一张门票包含多张饮品券，我们就用多个 \texttt{Coupon} 对象包裹该门票来装饰它。

\begin{codeboxt}{代码清单 14.29（程序 14.6 Enhanced-DecoratorDP）：\texttt{tester.cpp}}
#include "BaseTicket.h"
#include "Party.h"
#include "VIP.h"
#include "Coupon.h"

using namespace std;

void print(const string name, const Ticket * const ticket);

int main()
{
    Ticket *ticket = new BaseTicket();
    ticket = new Party(ticket);
    ticket = new VIP(ticket);
    ticket = new Coupon(ticket);    ①
    ticket = new Coupon(ticket);    ①

    print("John", ticket);
    delete ticket;

    ticket = new BaseTicket();
    ticket = new Coupon(ticket);
    ticket = new Coupon(ticket);
    ticket = new Party(ticket);
    ticket = new Coupon(ticket);

    print("Mary", ticket);
    delete ticket;

    ticket = new BaseTicket();
    print("Leslie", ticket);
    delete ticket;

    ticket = new BaseTicket();
    ticket = new Coupon(ticket);
    ticket = new VIP(ticket);

    print("Sidney", ticket);
    delete ticket;

    return 0;
}

void print(const string name, const Ticket * const ticket)
{
    cout << endl;
    cout << setw(26) << name + "'s ticket TOTAL: $"
         << ticket->get_cost() << endl;
    cout << "----------------------------" << endl;
}
\end{codeboxt}
\noindent{\fontsize{9bp}{12bp}\selectfont ① 多张饮品券}\par

结果与应用第一个版本相同，只是格式不同：

\begin{codebox}
      base ticket price: $30
    pregame party price: $25
      VIP seating price: $20
     drink coupon price: $ 5
     drink coupon price: $ 5

    John's ticket TOTAL: $85
----------------------------
      base ticket price: $30
     drink coupon price: $ 5
     drink coupon price: $ 5
    pregame party price: $25
     drink coupon price: $ 5

    Mary's ticket TOTAL: $70
----------------------------
      base ticket price: $30

  Leslie's ticket TOTAL: $30
----------------------------
      base ticket price: $30
     drink coupon price: $ 5
      VIP seating price: $20

  Sidney's ticket TOTAL: $55
----------------------------
\end{codebox}

\myGraphic{0.872}{images/CH14_F11_UN07_Mak.png}{}

仿照装饰器设计模式为门票应用建模所带来的好处包括

\begin{itemize}
\item \emph{动态增强的对象}——这是通过在运行时添加职责来实现的。
\item \emph{没有硬编码的增强项}——创建新的门票增强项（例如预留停车位）并不难。我们可以再添加一个 \texttt{Decorator} 子类，比如 \texttt{Parking}。移除我们不再需要的任何装饰器也会很容易。
\item \emph{松耦合且内聚的类}——各个门票装饰相互之间是松耦合的。每个装饰都只对自身的装饰负责。类 \texttt{Ticket} 与各个 \texttt{Decorator} 类之间也是松耦合的。门票无需知道自己是怎样被增强的（如果有的话）。
\item \emph{统一对待基础门票对象与已装饰门票对象}——我们不必为二者分别编写特殊代码。
\end{itemize}

\subsection{14.3.4 装饰器的通用模型}

图 14.12 展示了装饰器设计模式的通用模型。回想一下，正是依据设计模式的通用模型，我们才能针对某个架构问题创建定制的解决方案（8.1.3 节）。

\myGraphic{0.713}{images/CH14_F12_Mak.png}{图 14.12 装饰器设计模式的通用模型。请与图 14.11 对照。一个 \texttt{ConcreteComponent} 对象自身就是一个 \texttt{Component}。被 Decorator 对象包裹的 \texttt{ConcreteComponent} 同样也是一个 \texttt{Component}。超类 \texttt{Decorator} 的 component 成员变量要么指向一个 \texttt{ConcreteComponent} 对象，要么指向一个 \texttt{Decorator} 对象。这条链实现了 \texttt{Decorator} 对象的嵌套包裹。}

表 14.3 展示了示例应用如何应用该模式。

\begin{center}
{\bfseries 表 14.3 示例应用所应用的装饰器设计模式}\par\vspace{3bp}
\begin{tabularx}{\textwidth}{XX}
\toprule
设计模式 & 示例应用中的应用 \\
\midrule
超类 \texttt{Component} & 超类 \texttt{Ticket} \\
类 \texttt{ConcreteComponent} & 类 \texttt{BaseTicket} \\
超类 \texttt{Decorator} & 超类 \texttt{Decorator} \\
子类 \texttt{ConcreteDecoratorA}、\texttt{ConcreteDecoratorB} 等 & 类 \texttt{Party}、\texttt{VIP} 和 \texttt{Coupon} \\
成员函数 \texttt{operation()} & 成员函数 \texttt{get\_cost()} \\
\bottomrule
\end{tabularx}
\end{center}

\section*{小结}

\begin{itemize}
\item 单例设计模式为这样的代码建模：它必须确保在应用的运行期间，某个给定类的对象至多只能存在一个。我们通过把单例类的构造函数设为私有来控制该对象的创建。运行时由一个静态成员变量指向单例对象。
\item 单例设计模式还要求禁用单例类的拷贝构造函数和拷贝赋值运算符。这样可以防止在运行时创建单例对象的副本。我们必须谨慎，以免多线程应用创建出多个单例对象。
\item 组合设计模式为处理层级树结构数据的代码建模。
\item 组合设计模式为表示单个对象的类和表示组合对象的类提供相同的接口，从而简化代码。这使我们的代码能够统一地对待单个对象和组合对象。
\item 装饰器设计模式为这样的代码建模：它让我们能够在运行时扩展对象的职责。这是通过以属性和行为形式存在的装饰来实现的，我们可以把装饰实现为包裹器对象。
\item 装饰器设计模式减少了硬编码的需要，因为它使我们能够在不修改任何其他类的情况下，轻松地添加新装饰，或修改、移除已有的装饰。
\end{itemize}
